%% =============================================================================
%% Pre-thesis v24 — Chapter 5: Market Analysis and Feasibility  (FIRST HALF)
%% =============================================================================
%% This file is the direct continuation of Pre-thesis_v24_chapter4_part2.tex.
%% It must be included immediately after Chapter 4 in the master document.
%%
%% Sections covered in this file (first half):
%%   Chapter 5 opening paragraph
%%   Bangladesh Financial Landscape                            §5.1
%%   Underserved Population and Demand Analysis                §5.2
%%   Competitive and Comparative Analysis                      §5.3
%%   Revenue Projection                                        §5.4
%%     Interest spread mechanics, origination fee,
%%     deposit product yield
%%     [v24-F §7.1: revenue taxonomy table — NEW]
%%     [v24-F §7.1: no native-token hard-rule paragraph — NEW]
%%
%% v24 improvements applied in this file:
%%   * v24-F §7.1: Revenue stream taxonomy table inserted after
%%     the existing Revenue Projection text, before
%%     \subsection{Transaction Economics}.
%%     (Source: CWB_v24_Improvement_Analysis.md §6.1;
%%              crypto_exchange_ecosystem_dashboard analysis.)
%%   * v24-F §7.1: Hard-rule paragraph (no native CWB governance
%%     token in the prototype) added immediately after the taxonomy
%%     table, citing the FTX November 2022 collapse as evidence.
%%     (Source: CWB_v24_Improvement_Analysis.md §6.2.)
%%
%% Sections deferred to Pre-thesis_v24_chapter5_part2.tex:
%%   Transaction Economics                                     §5.4.1
%%   Cost Model                                                §5.4.2
%%   Break-even Analysis                                       §5.4.3
%%   Technical Feasibility                                     §5.5
%%   Operational Feasibility                                   §5.6
%%   Risk Analysis                                             §5.7
%%   Adoption and Scaling Pathway                              §5.8
%%   Bangladesh Regulatory Reality [v24-C §7.2]               §5.9
%% =============================================================================

% ─────────────────────────────────────────────────────────────────────────────
\chapter{Market Analysis and Feasibility}
\label{chap:market-analysis}
% ─────────────────────────────────────────────────────────────────────────────

The Crypto World Bank addresses a financing gap that conventional banking and
existing decentralised finance protocols have not closed: providing
collateral-flexible, algorithmically scored, and KYC-tiered credit to
economically active but financially excluded populations in emerging economies.
This chapter situates the CWB prototype within the Bangladesh financial
landscape, analyses the underserved target population, benchmarks the system
against comparable existing solutions, and evaluates the financial, technical,
and operational feasibility of the proposed four-tier architecture. The chapter
concludes with the Bangladesh regulatory reality and a phased adoption pathway
designed to build institutional trust before seeking a formal banking licence.
All market figures are presented at the 300-client, 6-bank prototype scale
unless otherwise stated; Section~\ref{sec:target-population} addresses the
fully-scaled addressable market separately.

% ─────────────────────────────────────────────────────────────────────────────
\section{Bangladesh Financial Landscape}
\label{sec:bangladesh-landscape}
% ─────────────────────────────────────────────────────────────────────────────

Bangladesh presents one of the most compelling deployment contexts for
blockchain-based financial inclusion infrastructure. With a population of
approximately 172~million and a nominal GDP of USD~460~billion in 2024~[R40],
the country has sustained above-6\% annual economic growth for over a decade,
driven primarily by the ready-made garments (RMG) sector, remittances
(USD~21.9~billion in FY~2024), and an expanding small and medium enterprise
(SME) ecosystem. Despite this macroeconomic trajectory, structural gaps in
financial access remain unresolved by the existing institutional architecture.

\begin{figure}[H]
\centering
\OnePageDiagram{fig-bangladesh-market.pdf}
\caption{Bangladesh financial landscape overview: formal banking penetration,
mobile financial services reach, microfinance sector coverage, and the
remaining unbanked population as the primary CWB target market. Population
figures from World Bank 2024~[R40]; MFS data from Bangladesh Bank Annual
Report~[R46].}
\label{fig:bangladesh-market}
\end{figure}

\paragraph{Formal banking exclusion.}
According to the World Bank Global Findex 2021 dataset~[R41], approximately
47\% of Bangladeshi adults lacked a formal bank account, placing the country
below the regional average for South Asia. Account ownership is concentrated
in urban and peri-urban households; in rural divisions --- Barisal, Rajshahi,
and Sylhet --- the unbanked share exceeds 55\%. Structural barriers include
minimum-balance requirements (typically BDT~2,500--5,000), branch geography
(the national bank-branch-to-population ratio is one branch per 17,500
people~[R41]), documentation demands that exclude informal workers without
payslips or land records, and the absence of collateral acceptable to
commercial banks from low-income borrowers. These barriers are not individual
failures; they are architectural properties of a credit system designed for
formal employment and registered property ownership.

\paragraph{Mobile financial services: reach without credit.}
The mobile financial services (MFS) layer has achieved remarkable reach while
solving only part of the financial access problem. bKash, operated by BRAC
Bank, surpassed 63~million registered wallets by Q3~2024~[R42]; Nagad,
operated by Bangladesh Post Office, reported 90~million registered
users~[R43]. However, both platforms remain principally peer-to-peer transfer
and utility-payment instruments. Neither extends credit to their user base;
they are velocity layers for existing liquidity, not sources of new capital.
The population that uses MFS daily --- small traders, garment workers,
agricultural households, rickshaw operators --- continues to rely on informal
moneylenders charging 25--60\% annual interest, rotating savings and credit
associations (ROSCAs, locally called \emph{Samity}), and microfinance
institutions whose onboarding processes can take 4--8~weeks~[R44].
CWB's AI-agent-assisted application path (Chapter~3,
Section~\ref{sec:aiml-support}) compresses this to hours.

\paragraph{Microfinance sector and its structural ceiling.}
Bangladesh hosts the world's most developed microfinance sector, including
BRAC (7.6~million active borrowers), Grameen Bank (founded 1983,
9.5~million members), and over 750~MFIs licensed by the Microcredit
Regulatory Authority (MRA)~[R44]. Aggregate MFI disbursements exceeded
BDT~1.8~trillion (approximately USD~16~billion) in FY~2024. However, MFI
penetration creates its own structural ceiling. Interest rates of 20--27\%
(flat, not declining balance) represent a high cost of capital for borrowers.
Group-meeting attendance requirements impose time costs incompatible with
the schedules of full-time workers. Most critically, credit-score portability
does not exist across MFI boundaries: a borrower who repays reliably at BRAC
for three years receives no formal credit benefit when seeking a loan from
Grameen Bank, a cooperative, or a commercial bank. CWB's on-chain Credit
Passport SBT (Chapter~3, Section~\ref{sec:credit-passport}) is designed to
solve precisely this portability problem --- every on-time repayment advances
the borrower's credit score permanently on-chain, regardless of which Local
Bank issued the loan.

\paragraph{SME credit gap.}
The International Finance Corporation estimates the SME finance gap for
Bangladesh at approximately USD~2.8~billion annually~[R45], representing
the difference between SME demand for formal credit and the supply delivered
by commercial banks and MFIs combined. Formal bank credit to SMEs is
constrained by collateral requirements --- typically land-registered property,
which excludes most urban micro-entrepreneurs --- and by documentation
processes averaging 6--10~weeks from application to disbursement. The CWB
four-tier architecture, with Local Banks operating within communities and the
AI agent compressing the application-to-decision cycle, is sized to serve
this segment at ticket sizes between 500~USDC and 5,000~USDC
(Table~\ref{tab:credit-tier-schedule} in Chapter~3).

\paragraph{Bangladesh Bank digital finance policy.}
Bangladesh Bank published its Digital Banking Regulation~2023~[R46] and the
updated FinTech Policy~2025~[R47], which together create a regulatory sandbox
framework for non-bank digital financial service providers. The sandbox allows
approved entities to operate for up to two~years under relaxed licensing
conditions while accumulating the track record required for full regulatory
authorisation under the Bank Company Act. CWB's adoption pathway
(Section~\ref{sec:adoption-pathway}) is designed around this sandbox
mechanism. The full regulatory pathway, including implications of the FATF
Travel Rule and anti-money-laundering requirements, is examined in
Section~\ref{sec:bangladesh-regulatory}.

% ─────────────────────────────────────────────────────────────────────────────
\section{Underserved Population and Demand Analysis}
\label{sec:target-population}
% ─────────────────────────────────────────────────────────────────────────────

The CWB architecture serves three economically distinct client segments, each
mapped to a KYC tier and Credit Passport tier range. The segmentation is
grounded in Bangladeshi income distribution data from the Household Income and
Expenditure Survey (HIES)~[R48] and qualitative field research on
microfinance demand published by the BRAC Institute of Governance and
Development~[R49].

\begin{table}[H]
\centering
\caption{CWB target population segments: income profile, primary financial
need, KYC tier alignment, and Credit Passport tier range.}
\label{tab:target-segments}
\small
\renewcommand{\arraystretch}{1.3}
\begin{tabular}{p{2.5cm}p{3.0cm}p{3.0cm}p{1.8cm}p{1.9cm}}
\hline
\textbf{Segment} & \textbf{Profile} & \textbf{Primary need} &
\textbf{KYC tier} & \textbf{Credit tier} \\
\hline
Urban micro-entrepreneur &
  Monthly income BDT~12,000--30,000; garment workers, street retail,
  rickshaw owners, domestic workers &
  Working-capital loan BDT~5,000--25,000 (50--250~USDC); short duration
  (3--6~months) &
  Tier~1 &
  Bronze--Silver \\
Rural agricultural household &
  Seasonal income BDT~6,000--18,000; smallholder farmers, pond-fish
  cultivators, poultry farmers &
  Seasonal credit BDT~5,000--50,000 (50--500~USDC); group lending preferred;
  repayment aligned to harvest cycle &
  Tier~1 &
  Bronze--Gold \\
SME / formal micro-business &
  Annual turnover BDT~0.5M--5M; tailoring workshops, small manufacturers,
  service providers, tech freelancers &
  Term loan BDT~50,000--500,000 (500--5,000~USDC); fixed deposit for
  surplus; longer duration (6--18~months) &
  Tier~2 &
  Silver--Platinum \\
\hline
\end{tabular}
\end{table}

\paragraph{Group lending demand.}
Bangladesh's MFI sector demonstrates robust demand for group-liability credit
structures. BRAC's core lending model is group-based; Grameen Bank has
disbursed over USD~38~billion through solidarity group lending since its
founding and reports a repayment rate above 97\%~[R44]. CWB's
\texttt{GroupLendingPool} (Chapter~3, Section~\ref{sec:group-lending})
replicates the group-liability structure on-chain, providing the same
social-collateral mechanism at a fraction of the administrative cost.
The five-member group structure directly maps to the \emph{Samity} group model
familiar to the vast majority of rural borrowers in Bangladesh, reducing
adoption friction for the primary target segment. The AI agent can coordinate
multi-member signature collection, track outstanding consents, and send
reminders to unsigned members without any Local Bank administrator
intervention (Chapter~4, Section~\ref{sec:aiml-support}).

\paragraph{Remittance corridor opportunity.}
Bangladesh receives the world's seventh-largest remittance inflow relative to
GDP. Workers in the Gulf Cooperation Council, Malaysia, and Singapore
collectively send USD~21.9~billion annually through formal channels, with an
estimated USD~6--8~billion flowing through \emph{hundi} (informal) channels
that carry lower fees but create compliance exposure for senders~[R48]. The
\emph{hundi} system persists because formal remittance corridors charge
3--7\% fees and take 2--5~business days for settlement. While remittance
corridors are scoped as Phase~3 future work (Chapter~6), the ERC-4337
wallet infrastructure built for CWB deposit accounts is directly compatible
with USDC remittance receipt, positioning CWB for corridor entry without
additional identity infrastructure investment.

\paragraph{Addressable market sizing.}
Conservatively restricting the addressable market to KYC~Tier~1 and
Tier~2 segments (Table~\ref{tab:target-segments}) across six of Bangladesh's
eight administrative divisions where Local Bank density is assumed to be
viable, the serviceable addressable market is approximately
12~million~households. At an average loan size of 150~USDC and 1.5~loan
cycles per household per year, the credit deployment potential reaches
USD~2.7~billion annually. At a 10\% net interest spread across the four-tier
hierarchy (Chapter~3 kinked rate model), this implies annualised protocol
net interest income of USD~270~million in a mature deployment, before
origination fee and fixed-deposit margin revenue. This ceiling is
illustrative; the prototype operates at the 300-client, 6-bank simulation
scale validated in the Foundry stress test (Chapter~4, US-3.15).

% ─────────────────────────────────────────────────────────────────────────────
\section{Competitive and Comparative Analysis}
\label{sec:competitive-analysis}
% ─────────────────────────────────────────────────────────────────────────────

CWB occupies a design space between traditional microfinance institutions,
decentralised finance (DeFi) lending protocols, and mobile financial services
platforms. No existing system in any of these categories offers the full
combination of: on-chain credit history portability, a four-tier capital
hierarchy, AI-agent-assisted banking operations, and locally-hosted
privacy-preserving inference. Table~\ref{tab:competitive-comparison}
benchmarks CWB against the most relevant comparators across eleven
architectural and regulatory dimensions.

\begin{table}[H]
\centering
\caption{CWB competitive and comparative benchmark across eleven dimensions.
\checkmark{} = natively supported; $\sim$ = partially or via third-party
integration; $\times$ = not supported; N/A = not applicable to platform
category.}
\label{tab:competitive-comparison}
\small
\renewcommand{\arraystretch}{1.3}
\begin{tabular}{p{3.5cm}p{1.5cm}p{1.5cm}p{1.7cm}p{1.5cm}p{1.7cm}}
\hline
\textbf{Feature} &
\textbf{CWB (v24)} &
\textbf{Aave v3} &
\textbf{Grameen} &
\textbf{bKash} &
\textbf{Goldfinch} \\
\hline
On-chain credit history portability & \checkmark & $\times$ & $\times$ & $\times$ & $\sim$ \\
Under-collateralised credit & \checkmark & $\times$ & \checkmark & $\times$ & \checkmark \\
AI-assisted loan approval & \checkmark & $\times$ & $\times$ & $\times$ & $\times$ \\
Agent-executed banking ops. & \checkmark & $\times$ & $\times$ & $\times$ & $\times$ \\
Tiered interest rate by borrower risk & \checkmark & $\sim$ & \checkmark & N/A & $\sim$ \\
Group lending pool & \checkmark & $\times$ & \checkmark & $\times$ & $\times$ \\
Privacy-preserving local inference & \checkmark & N/A & N/A & N/A & N/A \\
Reserve ratio on-chain transparency & \checkmark & \checkmark & $\times$ & $\times$ & $\sim$ \\
Chainlink Proof of Reserve & \checkmark & $\times$ & $\times$ & $\times$ & $\times$ \\
Mobile-first onboarding & \checkmark & $\sim$ & $\sim$ & \checkmark & $\times$ \\
FATF Travel Rule data packet & \checkmark${}^{*}$ & $\times$ & $\times$ & $\times$ & $\times$ \\
\hline
\multicolumn{6}{l}{\footnotesize ${}^{*}$ Specified; implementation is Future
Work (Section~\ref{sec:regulatory-compliance} and
Section~\ref{sec:bangladesh-regulatory}).} \\
\end{tabular}
\end{table}

\paragraph{DeFi protocol comparison.}
Aave~v3 and Compound~v3 are the leading collateralised DeFi lending
protocols, collectively managing over USD~15~billion in total value locked as
of Q1~2025. Both require overcollateralisation --- typically 150--200\% of
the loan value denominated in on-chain assets --- and have no concept of
borrower identity, credit history, or income. They are structurally
inaccessible to the CWB target population, who cannot post on-chain
collateral equivalent to, let alone larger than, the loan amount sought.
Goldfinch Finance and Centrifuge address real-world credit but operate via
institutional underwriters and are not designed for individual retail
borrowers in emerging markets; minimum deal sizes are USD~100,000--1,000,000.
CWB's SBT-based Credit Passport, AI-assisted ML scoring, and KYC-tier-gated
loan caps enable undercollateralised credit at ticket sizes starting at
50~USDC, which are three orders of magnitude smaller than any comparable DeFi
mechanism.

\paragraph{Traditional MFI comparison.}
BRAC and Grameen Bank are the most operationally successful microfinance
models in Bangladesh, and CWB explicitly models its group-lending structure
and institutional trust-bootstrapping pathway on BRAC's field
experience~[R49]. However, both institutions share structural limitations that
the CWB architecture addresses. Credit history does not transfer between MFI
institutions: a borrower who repays reliably at BRAC for three years has no
formal credit record visible to Grameen Bank or a commercial bank. Loan
processing takes 4--8~weeks from application to disbursement, driven by
manual document collection, physical approval meetings, and paper-based
disbursement processes. MFIs also do not offer savings products competitive
with commercial bank deposit rates; BRAC's savings rate of 4--5\% is below
the rate CWB's \texttt{FixedDeposit} product targets for Tier~2 depositors.

\paragraph{Mobile financial services comparison.}
bKash and Nagad occupy a complementary rather than competing position to CWB.
Neither platform extends credit; both are optimised for payment velocity at
low cost. CWB wallets are USDC-denominated EVM addresses rather than
MFS~wallet IDs, but the user experience pattern --- phone-number-linked,
SMS-verified, biometric-authenticated --- is designed to be immediately
familiar to the existing MFS user base, which numbered over 150~million
registered accounts across all operators as of 2024. The adoption pathway
(Section~\ref{sec:adoption-pathway}) describes a technical bridge between
CWB USDC balances and bKash wallets for on-ramp and off-ramp via authorised
exchange partners, allowing existing bKash users to fund CWB deposits without
purchasing cryptocurrency through an exchange.

\begin{figure}[H]
\centering
\OnePageDiagram{fig-competitive-positioning.pdf}
\caption{CWB competitive positioning map: two-axis plot of credit accessibility
(x-axis: collateral requirement from overcollateralised to zero-collateral)
versus institutional trust model (y-axis: centralised to decentralised).
CWB occupies the undercollateralised + decentralised-with-compliance quadrant,
which is unoccupied by any existing system in the Bangladesh market.}
\label{fig:competitive-positioning}
\end{figure}

% ─────────────────────────────────────────────────────────────────────────────
\section{Revenue Projection}
\label{sec:revenue-projection}
% ─────────────────────────────────────────────────────────────────────────────

The CWB revenue model is built around the four-tier interest spread as its
primary mechanism, with five supplementary revenue streams activated
progressively across deployment phases. Each tier in the hierarchy earns a
spread between the rate at which it borrows capital from the tier above and
the rate at which it lends to the tier below. The World Bank sets the base
rate published in the \texttt{interest\_rate\_tier} table; the National Bank
charges base~+~spread\textsubscript{NB} to Local Banks; the Local Bank
charges base~+~spread\textsubscript{NB}~+~spread\textsubscript{LB} to
clients; and the client's effective rate is then modulated downward by their
Credit Passport tier modifier, as specified in
Table~\ref{tab:credit-tier-schedule} of Chapter~3. The result is a system
where a Diamond-tier borrower (score 900--1000) pays a materially lower rate
than a Bronze-tier first-time borrower, creating a self-reinforcing incentive
for on-time repayment and credit-score improvement.

\paragraph{Interest spread mechanics.}
The kinked interest rate model (Chapter~3, Section~\ref{sec:oracle-architecture})
bounds the Local Bank lending rate below the kink utilisation threshold
$U^* = 80\%$ and applies a steeper incremental rate above it to discourage
pool over-utilisation. At the steady-state target of 70\% pool utilisation, a
Gold-tier borrower (credit score 550--749) pays the base rate minus 1.0\%
per annum. Five percent of all interest collected at each tier is diverted to
the \texttt{InsuranceFund} contract before upward distribution
(US-1.12,~Sprint~1), creating a reserve buffer against default losses. At
the 300-client prototype scale --- average outstanding balance of 500~USDC
per client, 70\% pool utilisation, 12\% lending rate --- gross annualised
interest income is approximately USD~18{,}000, and USD~17{,}100 after the
5\% insurance fund contribution. The \texttt{getReserveSummary()} view
function (US-1.13) reports these figures to the Reserve Transparency
Dashboard and to Chainlink Proof of Reserve in real time.

\paragraph{Loan origination fee.}
A flat origination fee of 0.5--1.0\% of the disbursed principal is deducted
at the point of disbursement by the \texttt{LoanController} contract. The fee
is retained by the Local Bank as an operating contribution, partially
offsetting the cost of the approver dashboard and the Chainlink Automation
registration. At an average loan size of 500~USDC and a 0.75\% fee, each
loan generates USD~3.75 in one-time origination revenue. Across 300~clients
executing 1.5~loan cycles per year at the prototype scale, annualised
origination revenue is approximately USD~1{,}688. At the fully scaled
12~million-household addressable market, the same fee rate applied to a
150~USDC average loan generates USD~270~million annually in origination
revenue alone, independent of the interest spread.

\paragraph{Deposit product margin.}
The \texttt{SavingsVault} pays a flexible deposit rate set by Local Bank
governance, bounded at or above the minimum defined in the
\texttt{interest\_rate\_tier} table to prevent predatory low-rate deposits.
The \texttt{FixedDeposit} product pays a term premium above the flexible
rate in exchange for a time-locked principal commitment. The spread between
the deposit rate paid to savers and the lending rate charged to borrowers
constitutes the net interest margin (NIM), the primary profit driver. With
flexible deposits at 4\% and lending at 12\%, the gross NIM before operating
costs and insurance fund contribution is 800~basis points per annum --- above
the 650~basis points reported by Grameen Bank in its FY~2024 financial
statements~[R44] and within the range of the most efficient commercial
microfinance operations globally.

% ── v24-F §7.1: Revenue taxonomy table ─────────────────────────────────────
%% NEW in v24: Revenue stream taxonomy table (CWB_v24_Improvement_Analysis.md §6.1)
%% Source: crypto_exchange_ecosystem_dashboard analysis + upgrade plan §7.1

\paragraph{Revenue stream taxonomy.}
Table~\ref{tab:revenue-streams} maps the six CWB revenue streams to their
institutional crypto exchange analogues, situating the CWB revenue model
within the broader landscape of digital financial infrastructure. The interest
spread is the primary mechanism and is the only revenue stream fully modelled
in the prototype; the remaining five streams are activated across deployment
phases as the client and institutional base grows.

\begin{table}[H]
\centering
\caption{CWB revenue stream taxonomy: analogues from the institutional crypto
exchange sector. Comparison drawn from ecosystem analysis of Binance, Coinbase,
Kraken, and KuCoin revenue structures~[R50]. Phase assignments follow the
adoption pathway (Section~\ref{sec:adoption-pathway}).}
\label{tab:revenue-streams}
\small
\renewcommand{\arraystretch}{1.3}
\begin{tabular}{p{3.5cm}p{3.2cm}p{5.2cm}}
\hline
\textbf{Revenue stream} & \textbf{Exchange analogue} & \textbf{CWB mechanism} \\
\hline
Interest spread (primary) &
  Binance OTC desk / institutional lending desk &
  \texttt{SavingsVault} deposit yield vs.\ retail lending rate; four-tier
  accumulated spread modulated downward by Credit Passport tier
  (Table~\ref{tab:credit-tier-schedule}). Fully modelled in the
  prototype. \\
Loan origination fee &
  Spot trading fee (Binance: 0.10\%; Coinbase: 0.40\%) &
  0.5--1\% flat fee deducted at disbursement by
  \texttt{LoanController}; retained by the Local Bank as an
  operating-cost contribution. Phase~1. \\
Deposit mobilisation &
  Binance Earn / KuCoin flexible savings &
  \texttt{SavingsVault} flexible deposits + \texttt{FixedDeposit}
  term products; Local Bank earns the SavingsVault--lending rate
  spread. Phase~1. \\
Local Bank registration &
  Exchange listing fee (\$500K+ per token listing) &
  One-time registration fee paid by a new Local Bank to the National
  Bank on institutional onboarding; covers \texttt{CREATE2} contract
  deployment gas and KYC verification of the institution. Phase~2. \\
Syndicated loan arrangement &
  Institutional OTC desk / block-trade facilitation &
  Lead Arranger fee on \texttt{TranchedPool} /
  \texttt{SyndicatedLoan} facilities for corporate borrowers above the
  KYC~Tier~2 cap; activated in Phase~2 for SME and cooperative
  borrowers. Phase~2. \\
Reserve transparency API &
  Exchange data-feed and analytics API subscriptions &
  Phase~3: a paid GraphQL API layer over The Graph subgraph,
  exposing reserve-ratio time series, pool utilisation, and insurance
  fund balance to institutional investors and regulators; public
  read-only tier free. Phase~3. \\
\hline
\end{tabular}
\end{table}

% ── v24-F §7.1: Hard rule — no native token ─────────────────────────────────
%% NEW in v24: No-native-token hard-rule paragraph (CWB_v24_Improvement_Analysis.md §6.2)
%% Source: FTX collapse analysis + upgrade plan §7.1

\paragraph{Hard rule: no native CWB governance token in the prototype.}
The FTX collapse of November 2022 demonstrated the systemic failure risk of
self-minted tokens used as collateral or as a primary reserve asset. FTX's
FTT token was simultaneously the dominant asset on Alameda Research's balance
sheet and the implicit collateral backstop for FTX's solvency; its value was
circular and collapsed to near zero within 72~hours of the Binance withdrawal
announcement, triggering an eight-billion-dollar solvency gap that wiped out
customer deposits across multiple jurisdictions. No native CWB governance
token is issued in the thesis prototype. The Credit Passport SBT provides
on-chain reputation, tier-gated credit access, and borrower identity without
creating a speculative asset whose value can be manipulated, used as circular
collateral, or become a systemic risk vector. A future governance token
(Phase~3, following institutional trust bootstrapping as described in
Section~\ref{sec:adoption-pathway}) is specified in Future Work
(Chapter~6, Section~\ref{sec:future-work}) with explicit anti-FTX collateral
hard rules encoded at the contract level: the token may not be accepted as
collateral within any tier of the CWB lending hierarchy; its issuance is
deferred until the system-wide reserve ratio has exceeded 150\% of
outstanding principal for six consecutive months; and its supply schedule
is governance-locked and cannot be inflated by any single keyholder or
multisig configuration below a 5-of-7 threshold.

%% =============================================================================
%% END OF CHAPTER 5 — FIRST HALF
%% Next: Pre-thesis_v24_chapter5_part2.tex
%% Continues from: \subsection{Transaction Economics}
%%
%% Sections remaining for Part 2:
%%   Transaction Economics          §5.4.1  (gas costs, oracle call costs)
%%   Cost Model                     §5.4.2  (infrastructure, audit, oracle)
%%   Break-even Analysis            §5.4.3  (300-client prototype scale)
%%   Technical Feasibility          §5.5
%%   Operational Feasibility        §5.6
%%   Risk Analysis                  §5.7    (table: 4 risk categories)
%%   Adoption and Scaling Pathway   §5.8    (BRAC pilot → Phase 1 → Phase 3)
%%   Bangladesh Regulatory Reality  §5.9    [v24-C §7.2: FATF para — NEW]
%% =============================================================================
